% §1 Introduction
Voice-based conversational interfaces are increasingly deployed in real-world applications,
from customer service \citep{openai2024realtimeapi} to personal assistants \citep{google2024geminivoice}.
Two competing architectural paradigms have emerged for building such systems.
\emph{Native voice models} process speech end-to-end within a single neural network,
eliminating inter-component boundaries and their associated latencies.
\emph{Cascade architectures} decompose the pipeline into three sequential components:
a speech recognition model that transcribes user audio to text,
a text-based language model that reasons and generates a response,
and a speech synthesis model that renders the response as audio.

Cascade architectures offer compelling practical advantages.
They are substantially cheaper to operate than native voice APIs \citep{openai2024realtimeapi}.
They allow the language model component to be replaced or updated independently.
They support tool use---the ability to call external APIs, databases, and services---through
standard text-based agent frameworks \citep{yao2023react,langgraph2024}.
They expose intermediate representations (transcripts, structured outputs) that aid
auditability, debugging, and safety monitoring.

Despite these advantages, cascade architectures suffer from a perceptual deficiency:
the \emph{silence gap}.
When a user finishes speaking, the cascade pipeline must complete speech recognition,
then wait for the language model to produce its first token,
then wait for the speech synthesis model to begin audio playback.
The sum of these delays produces an interval of silence that users perceive as unnatural.
Human conversation tolerates gaps of roughly 200 milliseconds before perceiving a response
as delayed \citep{levinson2015timing,stivers2009universals}.
Cascade pipelines routinely exceed this threshold by a factor of four to ten.

This gap is widely acknowledged in the engineering literature as a key limitation of cascade
voice systems, but it has received no formal treatment.
Practitioners address it through ad hoc engineering mitigations without a principled basis
for understanding which strategies reduce the gap structurally and which merely mask it
perceptually.
No prior work has formally characterized the irreducible latency floor of cascade voice pipelines as a function of pipeline topology, nor established formal conditions under which such systems can approach human conversational response norms.

This paper provides that framework.
We introduce the \emph{turn-taking threshold} $\tau_{human} = 200$~ms as the upper bound on
the \emph{silence gap}---the interval of silence between the end of the user's utterance and
the beginning of the agent's audio response---within which a \emph{cascade architecture} is
perceived as conversationally natural.
We make the following contributions:

\begin{enumerate}
\item \textbf{Formal pipeline model} (\S\ref{sec:pipeline_model}).
We represent the cascade voice pipeline as a directed acyclic graph of latency-contributing
stages and decompose each stage's latency into a \emph{reducible} component,
improvable with faster models or hardware, and an \emph{irreducible} component,
determined by architectural structure and eliminable only by changing the pipeline topology.

\item \textbf{Latency Floor Theorem} (\S\ref{sec:latency_floor}).
We prove that the silence gap of any cascade architecture has a lower bound $\Lambda^*$,
the sum of irreducible latencies across all pipeline stages,
that cannot be reduced by any model improvement or hardware upgrade.
We derive the \emph{Cascade Competitiveness Criterion} (CCC),
the formal condition under which a cascade architecture achieves perceptual equivalence
with a native voice model.

\item \textbf{ReAct extension} (\S\ref{sec:react_extension}).
We extend the formal model to tool-using agents built on the ReAct pattern \citep{yao2023react}.
We prove that each reasoning step adds a full latency cycle to the silence gap,
compounding the irreducible floor and making CCC satisfaction increasingly unlikely
for multi-step agents.

\item \textbf{Design lever taxonomy} (\S\ref{sec:design_levers}).
We classify six mitigation strategies by the latency component each targets,
the headroom each can recover, and its side effects.
We show that only two strategies---streaming synthesis and hybrid routing---reduce
the irreducible floor $\Lambda^*$ itself;
the remaining four are perceptual patches that mask the gap without reducing it.

\item \textbf{Open problems} (\S\ref{sec:open_problems}).
We identify four open research directions emerging from the formal analysis:
learned barge-in prediction, safety in cascade voice systems,
voice-aware agent planning, and principled hybrid routing policies.
\end{enumerate}

The remainder of this paper is organised as follows.
Section~\ref{sec:background} introduces the cascade pipeline architecture,
native voice models, and human turn-taking thresholds.
Section~\ref{sec:pipeline_model} develops the formal model.
Sections~\ref{sec:latency_floor} and~\ref{sec:react_extension} state and prove the main theorems.
Section~\ref{sec:design_levers} presents the design lever taxonomy.
Sections~\ref{sec:open_problems} and~\ref{sec:discussion} discuss implications.
Section~\ref{sec:conclusion} concludes.
